\subsection{A One-Sided Estimator}
\label{sec:design-accountant}

The estimator keeps three structures. A \emph{cumulative layer} is a pair of
Count-Min sketches, one over $\eps_{\mathrm{pure}}$ and one over $\rho$, so that
mixed pure-DP and zCDP charges compose without conversion on the hot path. An
\emph{epoch ring} is a bounded deque of epochs, each with its own pair of
sketches, supplying the windowed view of \eqref{eq:window}. A \emph{hot table}
holds the $H$ subjects with the largest accumulated charge exactly, because those
are the subjects whose budgets bind.

For key $s$ and cost $c$, the conservative update computes
$m_j=\min_{r} A_j[r,h_{j,r}(s)]$ per layer and raises every cell on $s$'s row
that falls below $m_j+c_j$ to $m_j+c_j$; the estimate is the minimum over rows.
The rule never decreases a cell and never increases an estimate, which is the
property we need. Standard Count-Min with plain increments also over-estimates,
but conservative update keeps the estimate closer to the truth for the same width
by refusing to charge collisions that the minimum already accounts
for~\cite{fusy2022phase,fusy2023countmin}. With $d=\lceil\ln(1/\delta_s)\rceil$
rows and $w=\lceil e/\eps_s\rceil$ columns, the estimate exceeds the true charge
by at most $\eps_s$ times the layer's total mass with probability at least
$1-\delta_s$, and the bound is one-sided: it bounds the excess from above rather
than bracketing an unbiased estimate.

The hot table interacts with the sketch in a way that needs care. An evicted
subject that returns must not restart from zero, because the sketch already holds
its past. We \emph{seed} it with the sketch's current upper bound and then
accumulate exactly. The estimate is therefore the sketch's bound on the past plus
the exact charge since re-admission, which still bounds the true total, with
excess at most the sketch's error at the moment of eviction. Seeding from below,
for instance with the point estimate minus one standard deviation, would restore
accuracy at the cost of soundness, and we do not do it.

The epoch ring is sized so that a window is always covered. With epochs of length
$\ell$ and a ring of $R$ epochs, a window straddling a boundary spans at most
$\lceil W/\ell\rceil+1$ epochs, so we require $R\ge\lceil W/\ell\rceil+1$ and set
the ring size to exactly that bound. An earlier version used
$R=\lceil W/\ell\rceil$, which silently dropped the oldest epoch of every
straddling window and made the windowed estimate an under-estimate by up to one
epoch of traffic. The fix is one line, but it is the difference between a
one-sided estimator and one that is not, and it is the kind of error a purely
cumulative test suite does not catch. Summing per-epoch bounds is itself a bound,
so the ring preserves one-sidedness without further argument.

Total state is $O(dw)$ for the cumulative layer, $O(Rdw)$ for the ring and $O(H)$
for the hot table, none of which depends on $N$. The default uses $w=2^{14}$, $d=4$,
$R=3$ and $H=1024$, occupying 4.7\,MB at $10^6$ subjects and not growing when the
population doubles. The price is an over-estimate that is invisible for heavy
hitters---the hot table holds them exactly---and bounded by $\eps_s$ times the layer
mass for the tail, where the cap is least likely to bind.
Section~\ref{sec:eval-accuracy} measures both effects.Section~\ref{sec:eval-accuracy} measures both
effects.
